% Generated draft: requires mathematical and source review.
\documentclass[11pt]{article}
\usepackage[margin=25mm]{geometry}
\usepackage{fontspec}
\setmainfont{Latin Modern Roman}
\usepackage{amsmath,amssymb,mathtools,unicode-math}
\setmathfont{Latin Modern Math}
\usepackage{enumitem,needspace,xurl,hyperref}
\setlength{\parindent}{0pt}
\setlength{\parskip}{5pt}
\setlength{\emergencystretch}{3em}
\newcommand{\N}{\mathbb{N}}
\newcommand{\Z}{\mathbb{Z}}
\newcommand{\Q}{\mathbb{Q}}
\newcommand{\R}{\mathbb{R}}
\newcommand{\C}{\mathbb{C}}
\newcommand{\F}{\mathbb{F}}
\newcommand{\A}{\mathbb{A}}
\newcommand{\PP}{\mathbb{P}}
\newcommand{\E}{\mathbb{E}}
\newcommand{\eps}{\varepsilon}
\newcommand{\dd}{\,\mathrm{d}}
\begin{document}
\section*{UnsolvedMath ID: 800386}
\textbf{Source problem code:} catalogue entry
\subsection{Definitions and mathematical statement}
Let \(K\) be a field of characteristic \(0\), let \(\bar K\) be an algebraic closure of \(K\), and let \(f:\mathbb P^N_K\to\mathbb P^N_K\) be a morphism defined over \(K\) of degree \(d\ge 2\). The source does not explicitly specify the permitted range of the projective dimension \(N\).

For a positive integer \(r\), define the M\"obius function by
\[
\mu(r)=\begin{cases}
1,&r=1,\\
(-1)^k,&r\text{ is a product of }k\text{ distinct primes},\\
0,&r\text{ is divisible by the square of a prime}.
\end{cases}
\]
The supplied extraction of Conjecture 9 states that there exists a positive integer \(n_0\), with no uniformity in the given data specified, such that for every integer \(n\ge n_0\),
\[
\#\operatorname{Per}_n^*\bigl(f,\mathbb P^N(\bar K)\bigr)
:=\sum_{\substack{m\mid n\\m\ge 1}}\mu\!\left(\frac{n}{m}\right)(d^m+1)^N.
\]
This reproduces the extracted relation, including its definition symbol \(:=\). The supplied record does not define \(\operatorname{Per}_n^*\), specify whether its count includes multiplicities, or explain whether the displayed relation is intended as a definition or a conjectural equality. Consequently, the intended mathematical assertion cannot be made unambiguous from this record alone. In particular, neither an interpretation as points of exact period \(n\) nor any additional assertion of positivity or nonemptiness is supplied here.
\subsection{Short English statement}
For a morphism of projective space of degree at least \(2\) over a field of characteristic \(0\), the extracted conjecture gives an eventual relation between an undefined starred periodic-point count and a M\"obius sum. The meaning of the starred count and the intended role of the definition symbol remain unresolved.
\subsection{Sources}
\begin{itemize}
\item Local UnsolvedMath registry, supplied source record: permanent ID \(800386\), exact source problem code \texttt{catalogue entry}; title: \emph{Open Problems in Arithmetic Dynamics --- Conjecture 9}. The registry labels the problem open; this is supplied status metadata, not an independently verified mathematical assessment.
\item Benedetto, DeMarco, Ingram, Jones et al., \emph{Current Trends and Open Problems in Arithmetic Dynamics} (2018), Conjecture 9. Supplied source URL: \url{https://arxiv.org/abs/1806.04980}. The registry attributes the proposal to these authors in 2018.
\item Supplied provenance records access on 2026-10-16 and extraction method \texttt{arxiv\_latex}. This access date and provenance are supplied curation information; no independent source visit or verification is claimed.
\end{itemize}
\end{document}
